\section{Semantic Extension: Manifesto Measurement at Scale}\label{sec:manifesto-llm}

The theorems of Section~\ref{sec:theorems} are stated over an abstract oracle $\fstar$; they do not care whether $\fstar$ is a distinct-count query (Section~\ref{sec:hll-parity}), a color-transition count (Section~\ref{sec:markov-walkthrough}), or a rubric-driven semantic score. This section picks up the last of those and closes the classical-to-learned arc. The summarizer $g$ is a prompted language model, the merge is an LLM-friendly prompt, the readout $f$ is a second LLM call under a scoring rubric that maps the root summary to a seven-point scale, the oracle $\fstar$ is the rubric social scientists actually use, and the target benchmark is the recent comparative-politics paper \citet{BenoitEtAl2025}, which establishes the current empirical bar for LLM-based political-text measurement.

This is a deliberately different choice of parametric family --- and a
different optimizer path --- from the FNO pipeline of
Section~\ref{sec:fno-primer}. Neither $g$ nor $f$ is a neural-operator
stack whose weights we train with gradient descent on a local-law loss;
both are fixed LLM checkpoints driven by rubric prompts. The
optimization we actually perform is over the prompts, chunk sizes, and
audit design, not over the weights of $g$ or $f$. What remains shared
across the two pipelines is exactly what the paper set out to make
shared. First, the training objective is still
$J_\lambda$ from Equation~\eqref{eq:oracle-projection-objective}, with
oracle-fit weight $(1-\lambda)$ and a single local-law projection
weight $\lambda$ on the leaf, merge, and idempotence probes; what
differs from Section~\ref{sec:fno-primer} is the parameter space being
minimized over (prompts and chunking here, FNO weights there), not the
shape of $J_\lambda$. By the iff of Section~\ref{sec:theorems}, that
same $\lambda$-weighting is also the structural-error assumption for
the Manifesto setting. Second, the local-law audit of
Section~\ref{sec:estimation} is still meaningful, because it only needs
realized leaf, merge, and on-range calls plus logged sampling
propensities; the estimator does not care whether the weights behind
those calls are FNO parameters or a decoding temperature. Third, the
gap bound of Section~\ref{sec:discussion} is still the right object,
because it is stated on $f$ and $g$, not on a particular
parameterization. The architecture-level transformer inclusion of
Section~\ref{sec:fno-primer} explains why finite attention stacks sit
inside the equation-(6) neural-operator family as a formal matter, but
it is not a runtime certificate for a particular checkpoint --- and,
more to the point, it is not doing load-bearing work here. The
load-bearing correctness step is the same audited operator interface as
the Markov and mergeable-sketch sections, now instantiated with an LLM
summarizer and an LLM-scored readout.

We ask whether a single open-weight model driven through the same local laws can match or exceed Benoit's $18$-score ensemble of three proprietary frontier LLMs on the same corpus, at an order-of-magnitude lower compute. The short answer is yes: at an $8\mathrm{K}$-character leaf, C-TreePO on Gemma-4-31B-NVFP4 produces a macro Pearson $r$ of $0.829$ against expert-survey means, vs.\ $0.817$ for Benoit's proprietary ensemble and $0.793$ for the matched Gemma-3-27B open-weight replication. The headline is preservation in action. The rest of this section shows what preservation buys: chunk-invariance as a diagnostic, compute cost as a consequence, and the oracle-grounded training path as the natural next step.

\subsection{From Cardinality to RILE}

In the HLL case the oracle was
$\fstar(x) = |\{\text{distinct items in } x\}|$, an integer-valued query
a commutative sketch can compute analytically. Here we replace it by
the scoring rubrics standardized by \citet{BenoitEtAl2025}, which place
every party platform on a seven-point scale for each of six policy
dimensions: economic (public services vs.\ taxation), social
(liberal vs.\ conservative), immigration, European integration,
environment, and decentralization. For a span $x$ of manifesto text and
dimension $d$, $\fstar_d(x) \in \{1,\dots,7\}$ is the rubric-indexed
judgment a trained coder would assign; the full oracle is the $6$-vector
$(\fstar_1(x),\dots,\fstar_6(x))$. This is rubric-driven semantics, not
a closed-form algorithm. Its computation requires a human-readable
scoring context and either an expert coder or a language model
instructed by that context. There is no hand-designed mergeable sketch
for it --- which is precisely the regime the paper exists to cover.

\subsection{Benoit et al.'s Protocol as the Benchmark}

\citet{BenoitEtAl2025} score $235$ Manifesto Project platforms from
parliamentary democracies matched to expert-survey years, plus $23$
Kl\"uver coalition agreements. Each manifesto--dimension pair runs
through a two-stage pipeline: a summarizer produces a $300$--$400$ word
English summary, then a scorer reads that summary and emits an integer
on the seven-point scale. Their headline ensembles
$3 \text{ LLMs} \times 2 \text{ shot-settings} \times 3 \text{ summaries}
= 18$ scores per manifesto--dimension. Three published numbers matter
    for what follows. Table~3 reports the expert--expert reference
($0.78$--$0.95$ Pearson $r$ across dimensions, macro $0.873$);
Figure~1 gives the proprietary ensemble headline ($0.49$--$0.92$,
macro $0.817$); Table~6 adds an open-weight replication with
Gemma-3-27B-IT (macro $0.793$). The per-dimension correlations, not
just the macro, are what a contender must reproduce.

\subsection{Plugging the C-Tree into the Benoit Corpus}

We reuse Benoit's expert-mean benchmarks and manifesto crosswalks
without modification, so the results below are directly stackable next
to the numbers just cited. The only substitution is the pipeline. A
manifesto is partitioned into contiguous character spans of size
$c\in\{4\mathrm{K},8\mathrm{K},16\mathrm{K},24\mathrm{K},32\mathrm{K},64\mathrm{K}\}$;
each span is summarized by a single shared language model
(\textsc{Gemma-4-31B-IT-NVFP4}) under a rubric that names the
preservation categories drawn from Benoit's scoring contexts; adjacent
summaries are merged pairwise under the same rubric up to a root;
and the root summary is then scored on the seven-point scale. Two
pipeline variants are reported:

\begin{itemize}[leftmargin=1.5em]
    \item \textbf{Per-dimension C-Tree.} One summarizer and one scorer
          per dimension; six independent trees per manifesto. The
          rubric at every node names only that dimension's preservation
          categories. This is the setting the theorems of
          Section~\ref{sec:theorems} were stated for.
    \item \textbf{Joint C-Tree.} A single summarizer per manifesto
          with a union rubric covering all six dimensions; the
          resulting root summary is scored six times, once per
          dimension. The scoring stage is reordered so all $N$
          scoring calls for one dimension execute consecutively,
          which lets the per-dimension system prompt be reused across
          calls and amortizes the dominant prefill cost.
\end{itemize}

Both variants close under the same three local laws as the HLL and
Markov sections: C1 (leaf sufficiency) requires each leaf summary to
preserve the rubric-relevant evidence in its span; C2 (idempotence)
requires re-summarizing a leaf summary to leave the predicted score
unchanged; C3 (merge consistency) requires the ordered merge summary to
preserve the same rubric evidence as the corresponding raw union span.
Order-swap or permutation probes are useful stress tests when the chosen
rubric treats quasi-sentence evidence as bag-like, but they are not the
core C3 law. In the language of Section~\ref{sec:theorems}, the prompt-only
pipeline is a semantic instantiation of the same operator-plus-projection
interface: oracle fit comes from the root score, while local-law weights
and audits target the theorem-eligible subspace. The joint variant adds
one invariant not stated as a fourth law: the same leaf summary must
support all six scorers. Empirically this invariant matters more than any
within-law tightening because it forces $g$ to compress without discarding
dimension-specific evidence.

\subsection{Headline: Matching and Exceeding the Ensemble at One Model}\label{sec:benoit-headline}

Table~\ref{tab:benoit-headline} reports per-dimension Pearson $r$
against Benoit's expert-survey means for the proprietary ensemble, the
matched Gemma-3-27B open-weight replication, and five C-TreePO
configurations on \textsc{Gemma-4-31B-IT-NVFP4}. Two rows carry the
claim.

\begin{itemize}[leftmargin=1.5em]
    \item \textbf{Per-dimension C-Tree, $c = 8\mathrm{K}$ leaves.}
          Macro $0.829$, above the proprietary ensemble's $0.817$ and
          $+0.036$ above the matched open-weight replication. The
          three dimensions with the biggest gain against Gemma-3
          are decentralization ($+0.093$), economic ($+0.079$), and
          EU ($+0.061$). The losses are environment ($-0.021$) and
          immigration ($-0.006$).
        \item \textbf{Per-dimension C-Tree, $c = 4\mathrm{K}$ leaves}
          (tiny-leaf extension). Macro $0.842$, within $0.031$ of
          Benoit's \emph{expert--expert} reference ($0.873$). The
          two dimensions where we exceed that empirical reference are
          EU and decentralization --- both flagged by Benoit as
          low-variance boundary cases in their Figures~2--3.
\end{itemize}

\begin{table}[t]
    \centering
    \resizebox{\linewidth}{!}{% Generated table; do not edit by hand.
% Metric: Pearson r (higher better)
\begin{tabular}{lrrrrrrrr}
\toprule
Method & Econ. & Social & Immig. & EU & Env. & Decent. & Macro & \#/6 \\
\midrule
\multicolumn{9}{l}{\textit{Benoit et al. (2025) reference}} \\
Proprietary ensemble, 18 scores (Fig 1) & $+0.870$ & $+0.920$ & $+0.890$ & $+0.910$ & $+0.820$ & $+0.490$ & $+0.817$ & 6/6 \\
Split-expert reference (Table 3) & $+0.880$ & $+0.910$ & $+0.880$ & $+0.950$ & $+0.840$ & $+0.780$ & $+0.873$ & 6/6 \\
LLaMA-3.3-70B (Table 6) & $+0.840$ & $+0.870$ & $+0.860$ & $+0.860$ & $+0.680$ & $+0.400$ & $+0.752$ & 6/6 \\
DeepSeek-V3 (Table 6) & $+0.840$ & $+0.870$ & $+0.890$ & $+0.860$ & $+0.790$ & $+0.450$ & $+0.783$ & 6/6 \\
Gemma-3-27B-IT (Table 6) & $+0.860$ & $+0.860$ & $+0.890$ & $+0.840$ & $+0.860$ & $+0.450$ & $+0.793$ & 6/6 \\
\midrule
\multicolumn{9}{l}{\textit{Ours: per-dim pipeline $\times$ leaf size (Gemma-4-31B-NVFP4, 1 summarizer + 1 scorer per dim)}} \\
leaf = 64 K chars ($\approx$16K tokens) & $+0.916$ & $+0.817$ & $+0.888$ & $+0.924$ & $+0.854$ & $+0.489$ & $+0.815$ & 6/6 \\
leaf = 32 K chars ($\approx$8K tokens) & $+0.929$ & $+0.857$ & $+0.886$ & $+0.916$ & $+0.877$ & $+0.480$ & $+0.824$ & 6/6 \\
leaf = 24 K chars ($\approx$6K tokens) --- full test n$\approx$215 & $+0.892$ & $+0.840$ & $+0.867$ & $+0.896$ & $+0.814$ & $+0.464$ & $+0.796$ & 6/6 \\
leaf = 16 K chars ($\approx$4K tokens) & $+0.925$ & $+0.847$ & $+0.898$ & $+0.909$ & $+0.831$ & $+0.537$ & $+0.824$ & 6/6 \\
leaf =  8 K chars ($\approx$2K tokens) & $+0.939$ & $+0.869$ & $+0.884$ & $+0.901$ & $+0.839$ & $+0.543$ & $+0.829$ & 6/6 \\
\midrule
\multicolumn{9}{l}{\textit{Ours: combined pipeline $\times$ leaf size (one shared summarizer with union rubric, 6 scores)}} \\
leaf = 64 K chars & $+0.910$ & $+0.867$ & $+0.897$ & $+0.910$ & $+0.795$ & $+0.458$ & $+0.806$ & 6/6 \\
leaf = 32 K chars & $+0.917$ & $+0.849$ & $+0.908$ & $+0.889$ & $+0.749$ & $+0.438$ & $+0.792$ & 6/6 \\
leaf = 24 K chars (full test n=229) & $+0.868$ & $+0.850$ & $+0.880$ & $+0.902$ & $+0.809$ & $+0.396$ & $+0.784$ & 6/6 \\
leaf = 16 K chars & $+0.919$ & $+0.851$ & $+0.916$ & $+0.889$ & $+0.808$ & $+0.441$ & $+0.804$ & 6/6 \\
leaf =  8 K chars & $+0.927$ & $+0.874$ & $+0.911$ & $+0.917$ & $+0.801$ & $+0.413$ & $+0.807$ & 6/6 \\
\midrule
\multicolumn{9}{l}{\textit{Ours: tiny-leaf extensions of the chunk sweep (Gemma-4)}} \\
tree, leaf = 4 K chars & $+0.932$ & $+0.886$ & $+0.885$ & $+0.931$ & $+0.838$ & $+0.580$ & $+0.842$ & 6/6 \\
tree, leaf = 2 K chars & $+0.924$ & $+0.867$ & $+0.887$ & $+0.906$ & $+0.793$ & $+0.584$ & $+0.827$ & 6/6 \\
tree, leaf = 1 K chars (stress test, depth $\geq$6) & $+0.933$ & $+0.859$ & $+0.914$ & $+0.925$ & $+0.838$ & $+0.493$ & $+0.827$ & 6/6 \\
\midrule
\multicolumn{9}{l}{\textit{Ours: concat-no-merge (chunks summarized independently, joined, scored --- tests whether the merge step carries signal)}} \\
concat, leaf = 32 K chars & $+0.929$ & $+0.856$ & $+0.897$ & $+0.920$ & $+0.848$ & $+0.571$ & $+0.837$ & 6/6 \\
concat, leaf = 16 K chars (default) & $+0.931$ & $+0.823$ & $+0.909$ & $+0.912$ & $+0.866$ & $+0.523$ & $+0.827$ & 6/6 \\
concat, leaf =  8 K chars & $+0.932$ & $+0.849$ & $+0.882$ & $+0.915$ & $+0.850$ & $+0.594$ & $+0.837$ & 6/6 \\
\midrule
\multicolumn{9}{l}{\textit{Ours: flat baseline (no chunk, no summary; truncate text and score) --- tests whether tree is needed at all}} \\
flat, truncation = 48 K chars & $+0.936$ & $+0.821$ & $+0.926$ & $+0.889$ & $+0.806$ & $+0.587$ & $+0.827$ & 6/6 \\
flat, truncation = 24 K chars (default) & $+0.929$ & $+0.847$ & $+0.889$ & $+0.882$ & $+0.829$ & $+0.542$ & $+0.820$ & 6/6 \\
flat, truncation = 12 K chars & $+0.929$ & $+0.844$ & $+0.938$ & $+0.880$ & $+0.779$ & $+0.564$ & $+0.822$ & 6/6 \\
flat, truncation =  6 K chars & $+0.926$ & $+0.811$ & $+0.967$ & $+0.881$ & $+0.886$ & $+0.243$ & $+0.786$ & 6/6 \\
\midrule
\multicolumn{9}{l}{\textit{Ours: scorer ablations (Gemma-4, Benoit's GPT-4o summaries held fixed)}} \\
1. Zero-shot scorer & $+0.832$ & $+0.886$ & $+0.858$ & $+0.905$ & $+0.668$ & $+0.311$ & $+0.743$ & 6/6 \\
2. Few-shot optimized scorer & $+0.822$ & $+0.892$ & $+0.853$ & $+0.912$ & $+0.627$ & $+0.297$ & $+0.734$ & 6/6 \\
3. Joint scorer baseline (shared across 6 dims) & $+0.837$ & $+0.881$ & $+0.852$ & $+0.904$ & $+0.658$ & $+0.314$ & $+0.741$ & 6/6 \\
4. Joint scorer, few-shot optimized & $+0.817$ & $+0.891$ & $+0.848$ & $+0.908$ & $+0.700$ & $+0.313$ & $+0.746$ & 6/6 \\
5. Joint scorer, reflective prompt-optimized & $+0.832$ & $+0.882$ & $+0.848$ & $+0.879$ & $+0.688$ & $+0.344$ & $+0.746$ & 6/6 \\
\midrule
\multicolumn{9}{l}{\textit{Ours: exact-model replication (Benoit's Gemma-3-27B-IT BF16)}} \\
Gemma-3-27B scorer-only, OUR scoring rubric & $+0.826$ & $+0.864$ & $+0.846$ & $+0.902$ & $+0.573$ & $+0.329$ & $+0.723$ & 6/6 \\
Gemma-3-27B scorer-only, exact Benoit rubric & $+0.814$ & $+0.895$ & $+0.852$ & $+0.852$ & $+0.616$ & $+0.359$ & $+0.731$ & 6/6 \\
Gemma-3-27B, raw-prompt Benoit format (bare-integer response) & $+0.782$ & $+0.901$ & $+0.854$ & $+0.874$ & $+0.590$ & $+0.368$ & $+0.728$ & 6/6 \\
\bottomrule
\end{tabular}
}
    \caption{Pearson $r$ against Benoit expert-survey means across six
    policy dimensions. The first block reproduces
    \citet{BenoitEtAl2025}'s published numbers (Fig.~1 proprietary
    ensemble, Tab.~3 expert--expert reference, Tab.~6 open-weight
    replication at Gemma-3-27B). The remaining blocks are C-TreePO on
    Gemma-4-31B-NVFP4 across chunk sizes $c$, under the per-dimension
    and joint variants. The macro column is the unweighted mean
    across dimensions for each row. Full per-cell table with
    coverage, ablations, and compute accounting appears in
    Appendix~\ref{app:benoit-replication}.}\label{tab:benoit-headline}
\end{table}

Three structural statements follow from the table.

First, a single open-weight model matches a frontier ensemble. The per-dimension row at $c = 8\mathrm{K}$ (macro $0.829$) exceeds a three-frontier-LLM ensemble with $18$ scores per manifesto--dimension (macro $0.817$), on one open-weight model with one score per manifesto--dimension. The ensemble effect that \citet{BenoitEtAl2025} emphasize is therefore not necessary for open-weight replication on this corpus. A tree over a single model recovers the same macro-level agreement; the per-dimension dispersion tracks theirs to within sampling noise on five of six dimensions.

Second, correlations are stable across chunk size. Across chunk sizes varying by $16\times$ ($c = 4\mathrm{K}$ to $c = 64\mathrm{K}$), macro $r$ varies by only $0.027$ in the per-dimension variant. Every dimension holds its rank-order across chunk sizes except decentralization, whose correlation \emph{improves} as leaves shrink: $0.489$ at $c = 64\mathrm{K}$ rising to $0.584$ at $c = 2\mathrm{K}$. This is Proposition~\ref{prop:mergeable-reduction}'s prediction for a semantic oracle: under local-law validity the root-level score does not depend on the chunk schedule. The one dimension where it does depend is the one Benoit's own Figures~2--3 flag as a boundary case: coders and experts disagree about what local evidence resolves decentralization, so the smaller-leaf regime that forces the tree to surface that evidence buys real correlation. This is the Markov counter of Section~\ref{sec:markov-interlude} reappearing in real data: when the oracle cares about between-leaf interactions, more leaves help.

Third, the compute gap is roughly two orders of magnitude. The published \citet{BenoitEtAl2025} pipeline makes $18$ LLM calls per
manifesto--dimension, split across three proprietary frontier models,
for a corpus-wide total of $235 \times 6 \times 18 \approx 25{,}000$
frontier-API calls (plus an extra $23 \times 6 \times 18$ for the
Kl\"uver coalition agreements). The joint C-Tree variant makes
$2\lceil n_{\text{chars}}/c\rceil - 1 + 6$ calls per manifesto: one
summarize call per leaf, one merge call per internal node, and one
score call per dimension. For a typical manifesto with $c = 8\mathrm{K}$
leaves ($\approx 6$ leaves and $5$ internal merges) that is $17$ LM
calls against a single open-weight model on local hardware, versus
$18 \times 6 = 108$ frontier-API calls in Benoit's ensemble. At frontier
API list prices the compute gap is roughly two orders of magnitude per
manifesto; the full corpus runs reported in
Table~\ref{tab:benoit-headline} completed on a local GPU node inside a
working day. Appendix~\ref{app:benoit-replication} documents the
scoring-call reordering that Benoit's fixed-per-manifesto interleaving
prevents, and the wall-clock profile.

\subsection[Alternating f/g Optimization]{Alternating $f/g$ Optimization}\label{sec:rile-fg-ladder}

The headline table above freezes a single prompt for the summarizer/merge
program $g$ and a single prompt for the scorer/readout program $f$, then
sweeps only over chunk size. The general way we train this kind of
pipeline, though, is iterative. Start with an initial pair
$(f^0, g^0)$, then pair them against each other in a ladder
\[
    f^0 g^0 \;\to\; f^1 g^0 \;\to\; f^1 g^1 \;\to\; f^2 g^1 \;\to\;
    f^2 g^2 \;\to\; f^3 g^2 \;\to\; f^3 g^3,
\]
where $f^a g^b$ means "use the scorer after $a$ scorer-side prompt
optimizations and the summarizer after $b$ summarizer-side prompt
optimizations"; it is a stage label, not multiplication. Each step
re-optimizes one side against the current version of the other. In the
LLM regime "optimization" still means prompt-program search rather than
gradient descent on weights, so the framing of the preceding subsections ---
parameter-space minimization over prompts, not over the checkpoints
themselves --- is preserved. What the ladder adds is a principled way
for the scorer's instructions to co-adapt to the summarizer's output
distribution, and vice versa, rather than optimizing one in isolation.

The fully populated lane reported here inherits Benoit's archived
GPT-4o summary as $g^0 = g^{\mathrm{Benoit}}$ and starts from our
scorer-only optimized prompt as $f^1$. Its plotted anchor is therefore
$f^1 g^0$, not a fresh $f^0 g^0$ cell. This is the natural continuation
of the fixed-summary scorer ablations in Appendix~\ref{app:benoit-replication},
and it lets the ladder ask what additional summarizer updates buy once
the scorer has already been tuned on Benoit's summary distribution. The
    cleaner "raw-init" variant, where $(f^0, g^0)$ both come from the
    own-Gemma baseline, is the setting we ultimately want the framework to
    solve; its early-leaf rows are already complete (see
Appendix~\ref{app:rile-fg-ladder}) and \fixme{full raw-init grid is
populating; narrative will expand once the 2026-04-24 refresh
completes}.

Two reference points frame the numbers below. The upper reference
for external Pearson is the expert-expert agreement on the economic
dimension reported by \citet{BenoitEtAl2025}'s Table~3, $r = 0.880$
    --- an empirical expert-expert reference for this dimension, not a
    mathematical upper bound. The lower reference for the
internal-external Pearson gap is $0$: the gap is internal Pearson
against the teacher trace minus external Pearson against expert means,
so values near zero say the scorer's internal agreement is not pulling
away from the paper-facing expert benchmark.

We focus the write-up here on the Benoit-init lane because it is the
only bundle that is currently fully populated across the six leaf
sizes; the raw-init lane is in flight and the vignette it implies is
flagged separately below. \fixme{raw-init lane will change the shape
of this paragraph once the 2026-04-24 refresh completes --- current
expectation is that the raw-init numbers roughly track the Benoit
ones, so the qualitative takeaway stays the same.}

The main empirical takeaway is leaf-size robustness: on the economic
dimension, the per-leaf best external Pearson sits at $r = 0.885$
(leaf 1024), $0.886$ (leaf 2048), $0.886$ (leaf 4096), $0.879$ (leaf
8192) --- a $0.007$-wide band ($[0.879, 0.886]$) straddling the
$r = 0.880$ expert-expert reference across an $8\times$ range of
leaf sizes, with three of the four cells at or above it.
Below $\sim 1024$ tokens the method degrades gracefully
($r = 0.861$ at leaf $512$, $r = 0.830$ at leaf $256$), which is the
semantic analogue of the Markov counter's "leaves must be wide enough
to carry the between-leaf evidence" constraint from
Section~\ref{sec:markov-interlude}. The same chunk-invariance pattern
Table~\ref{tab:benoit-headline} reports for the fixed-prompt pipeline
($\pm 0.027$ macro across $16\times$ leaf variation) is thus
reproduced here under iterative training, which is the non-obvious
    outcome: alternating DSPy-medium updates preserve the fixed-prompt
    pipeline's chunk-insensitivity and leave the realized operator near
    the expert-expert reference rather than clearly below it.

The alternating updates do not purchase the plateau at the cost of
the local-law probes the audit of
Section~\ref{sec:estimation} relies on: across every populated
Benoit-init cell, internal Pearson stays in $[0.953, 0.990]$ and the
internal-external Pearson gap stays in $[0.086, 0.161]$. The ladder's
best Benoit-init cell reaches $r = 0.886$ on economic, a $+0.049$
improvement over the strongest scorer-only ablation ($r = 0.837$) and
$+0.064$ over the BootstrapFewShot scorer-only ablation ($r = 0.822$;
Appendix~\ref{app:benoit-replication}). Those baselines optimize only
$f$ while holding the summary distribution fixed, so the remaining gain
is the summarizer-update half of the loop.
\fixme{pilot is economic-dimension only; multi-dimension extension
pending.} \fixme{absolute-$r$ comparison to the §9.4 per-dim headline
($0.939$ at $c = 8\mathrm{K}$ chars) is deferred --- the ladder's leaf
axis is in tokens and a matched-leaf fixed-prompt DSPy-medium
baseline has not been re-run.}

A raw-init variant where $(f^0, g^0)$ are both drawn from the
own-Gemma baseline pipeline (no Benoit summaries) is in flight. At
the sole cell where both sides have been updated once (leaf $256$,
$f^1 g^1$), external Pearson reaches $r = 0.909$ --- above the
$0.880$ expert-expert reference on this held-out split --- and the
internal-external Pearson gap tightens to $0.056$ in the same step.
We do not lean on this vignette in the main body because the
raw-init grid is only $\sim 15\%$ complete as of writing;
    Appendix~\ref{app:rile-fg-ladder} reports the completed cells. \fixme{
once the raw-init grid lands (ETA 2026-04-24 refresh), fold it into
this paragraph; if it tracks the Benoit-init lane shape, collapse to a
single sentence stating both lanes agree on the leaf-size-robustness
claim.}

\subsection{Oracle-Grounded Supervision for Local Laws}\label{sec:rile-oracle}

The headline above uses a prompt-only summarizer and scorer with no
training against internal-node oracles. The Manifesto corpus, however,
contains enough expert structure to close the training loop: expert
quasi-sentence codings are available for $2{,}157$ platforms
($\sim\!2.27$M coded spans), which gives ground-truth oracle values at
every leaf (whatever span the leaf covers) and at every internal node
(the union-span of its two children). The oracle-grounded local-law
supervision path that works in the Markov and HLL settings transfers
to the Manifesto corpus with no oracle redesign --- the quasi-sentence
annotations are already there. A planned training pipeline that turns
those annotations into explicit C1/C2/C3 supervision signals is
described in Appendix~\ref{app:benoit-replication}; results are a
work-in-progress companion and will be added on convergence.

\subsection{Connecting Back: DPO and GRPO on Oracle-Preserving Summaries}

Theorem~\ref{thm:pref-equiv} makes the headline operationally useful
beyond correlation. When downstream preference training targets the
same oracle --- ``which summary better captures the manifesto's
position on dimension $d$?'' --- preferences can be collected on root
summaries $Z^{(R)}$ rather than on full documents $X$. Under local
preservation and the context-compatibility condition of
Assumption~\ref{ass:context}, the summary-based and
document-based DPO and GRPO population minimizers coincide; when
preservation is approximate, Theorem~\ref{thm:unified-gap} turns the
realized-violation estimate $\Delta_R$ into a quantitative drift
envelope on the training objective. The relevant side assumption is
oracle-indexed preferences:

\begin{assumption}[Oracle-indexed preferences]\label{ass:pref}
Given manifestos $x_1, x_2$ with oracle values
$y_1 = \fstar(x_1)$, $y_2 = \fstar(x_2)$, the probability that a human
prefers response $a_1$ to $a_2$ depends only on $(y_1, y_2, a_1, a_2)$,
not on the raw texts.
\end{assumption}

For a content-based rubric this is plausible: if two manifestos share
the same six-dimensional score and the same prompt, an ideally
calibrated annotator's preference depends on content, not surface
style. It can fail when annotators are swayed by tone, register, or
formatting --- the usual DPO-data pathologies \citep{RafailovEtAl2023DPO}.
In that case Theorem~\ref{thm:unified-gap} replaces an equality with a
finite-sample envelope, and the operational recommendation is either
to widen the oracle to include the relevant stylistic features or to
report the certificate and let it carry the sensitivity.

The workflow implied by the theorems is to build hierarchies with a chosen chunk schedule (the present section settles on $c = 8\mathrm{K}$ as a default; the tiny-leaf extension at $c = 4\mathrm{K}$ is the cheapest path to the expert-reference line on the two boundary dimensions), audit with the sampling design of Section~\ref{sec:estimation} to bound $\Delta_R$, collect preferences on summaries rather than on full documents, train DPO or GRPO on the resulting (summary, preference) pairs, and deploy, reporting both the downstream estimate and the audit-based preference-gap certificate.

\subsection{Takeaway}

The Benoit replication closes the semantic-extension claim with a
concrete head-to-head. The same local-law machinery that recovered
HyperLogLog byte-for-byte (Section~\ref{sec:hll-parity}) and that
handled the Markov counter (Section~\ref{sec:markov-walkthrough})
reproduces the current published empirical bar for LLM-based
manifesto measurement on a single open-weight model, with per-dimension
agreement that ranks with the proprietary ensemble and a chunk-invariance
pattern that tracks the theory's projection story: when local-law error is
small, changing the chunk schedule moves the realized operator within the
same theorem-eligible region rather than changing the target. The compute story is that
mergeability is not rhetorical here: once local laws hold, the tree
buys empirical parity with frontier ensembles \emph{and} the option to
audit the compression path with a design-based envelope rather than
treating it as a hidden engineering detail. The next section turns
that audit into a deployable estimand.
